\documentclass[11pt,a4paper]{article}

% ---------- Packages ----------
\usepackage[margin=2.2cm]{geometry}
\usepackage[utf8]{inputenc}
\usepackage[T1]{fontenc}
\usepackage{amsmath}
\usepackage{booktabs}
\usepackage{longtable}
\usepackage{array}
\usepackage{tabularx}
\usepackage{ragged2e}
\newcolumntype{Y}{>{\RaggedRight\arraybackslash}X}
\usepackage{xcolor}
\usepackage{enumitem}
\usepackage{listings}
\usepackage{parskip}
\usepackage{titlesec}
\usepackage{fancyhdr}
\usepackage[hidelinks,bookmarksopen=true,bookmarksnumbered=true]{hyperref}

% ---------- Colors ----------
\definecolor{codebg}{RGB}{245,245,245}
\definecolor{codeframe}{RGB}{200,200,200}
\definecolor{headingblue}{RGB}{20,60,110}
\definecolor{mustcolor}{RGB}{150,20,20}
\definecolor{donecolor}{RGB}{20,110,40}

% ---------- Section styling ----------
\titleformat{\section}{\large\bfseries\color{headingblue}}{\thesection}{0.6em}{}
\titleformat{\subsection}{\normalsize\bfseries\color{headingblue}}{\thesubsection}{0.6em}{}
\titleformat{\subsubsection}{\normalsize\bfseries}{\thesubsubsection}{0.6em}{}

% ---------- Code listing style ----------
\lstset{
  basicstyle=\ttfamily\small,
  backgroundcolor=\color{codebg},
  frame=single,
  rulecolor=\color{codeframe},
  breaklines=true,
  columns=fullflexible,
  showstringspaces=false,
  xleftmargin=4pt, xrightmargin=4pt,
  aboveskip=6pt, belowskip=6pt
}

\newcommand{\code}[1]{\texttt{\small #1}}
\newcommand{\done}{\textcolor{donecolor}{\textbf{DONE}}}
\newcommand{\pending}{\textcolor{mustcolor}{\textbf{NOT STARTED}}}

\setlength{\headheight}{14.5pt}
\addtolength{\topmargin}{-2.5pt}
\pagestyle{fancy}
\fancyhf{}
\fancyhead[L]{QuantumDx --- Project Context}
\fancyhead[R]{\thepage}
\renewcommand{\headrulewidth}{0.4pt}

% Compact "phase status line" macro: renders as a small aligned block
% instead of one long \quad-separated line, avoiding margin overflow.
\newcommand{\phasestatus}[5]{%
\begin{quote}\small\raggedright
\textbf{Status:} #1 \\
\textbf{Code:} #2 \\
\textbf{Docs:} #3 \\
\textbf{Tests:} #4 \\
\textbf{Results:} #5
\end{quote}}

\title{\textbf{QuantumDx --- Complete Project Context}\\[0.3em]
\large Hybrid Quantum Machine Learning Platform for Early Cardiovascular Disease Detection\\[0.2em]
\normalsize Smart India Hackathon 2026 --- Problem Statement SIH26139}
\author{Handoff document for continuing project work across AI sessions}
\date{Generated: Phase 5 complete, awaiting Phase 6 (September 2026)}

\sloppy
\begin{document}

\maketitle
\thispagestyle{empty}

\begin{abstract}
\noindent This document is a complete, self-contained record of the QuantumDx project: what problem it solves, every design decision made, every phase completed so far (Phase 1 through Phase 5), the exact numerical results obtained, the engineering principles enforced throughout, the current repository structure, and the remaining phases still to be built. It is written so that a reader (human or AI assistant) with \emph{zero prior context} can read it once and be immediately productive continuing the project --- verifying claims against the repository, running the next phase, or reasoning about design trade-offs already settled. Every number quoted here comes from an actual executed run, not an estimate, unless explicitly marked otherwise.
\end{abstract}

\tableofcontents
\newpage

% =========================================================================
\section{Project Overview}
% =========================================================================

\subsection{The problem statement}

\textbf{SIH26139 --- Hybrid Quantum Machine Learning Platform for Early Disease Detection.} The problem statement (Smart India Hackathon 2026) asks for a platform that integrates classical preprocessing and feature engineering with quantum-enhanced learning models (QSVM, QNN, or VQC), applied to biomedical data, supporting data ingestion, hybrid model training, prediction, explainability, and performance evaluation against purely classical baselines. Central requirements: a hybrid quantum-classical architecture, classical preprocessing, feature engineering/selection, quantum-enhanced ML, prediction, explainability, performance evaluation, comparison with classical baselines (accuracy, sensitivity, specificity), computational efficiency, generalization, scalability, quantum simulator compatibility, near-term quantum hardware compatibility, and comprehensive documentation.

\subsection{Our product decision}

\begin{itemize}[leftmargin=1.4em]
  \item \textbf{Disease:} Cardiovascular disease.
  \item \textbf{Initial task:} Early cardiovascular disease risk detection (binary classification).
  \item \textbf{Development dataset:} UCI Heart Disease dataset (Cleveland subset), used explicitly as a \emph{development/benchmark} dataset, not a claim about high-dimensional biomedical data in general. The architecture is deliberately designed so a larger cardiovascular dataset can be substituted later via configuration, not a rewrite.
  \item \textbf{Project name:} QuantumDx.
\end{itemize}

\subsection{The governing documents}

Two large planning documents were produced before any code was written, and both remain the source of truth for scope and philosophy:

\begin{itemize}[leftmargin=1.4em]
  \item \textbf{\code{PRD.md}} --- the full Product Requirements Document. 43 sections covering executive summary, problem statement, product goals/non-goals, personas, functional and non-functional requirements, the complete ML pipeline, classical and quantum requirements, explainability, benchmarking, robustness, dashboard/API/architecture, a SIH26139 requirement-to-feature traceability matrix, and acceptance criteria. \textbf{Never modified after creation.}
  \item \textbf{\code{MVP\_SPEC.md}} --- the MVP Specification and Implementation Plan. Defines the \emph{smallest complete} system satisfying SIH26139: the exact dataset schema, the 13-phase development plan (Phase 0 through Phase 12), project structure, engineering practices, and a real-hardware validation plan (Section~8.5, described in \S\ref{sec:ibm-plan} below). \textbf{Never modified after creation.}
\end{itemize}

Both documents establish, and every phase since has upheld, three non-negotiable rules:
\begin{enumerate}[leftmargin=1.4em]
  \item \textbf{Depth loses to completeness.} A working end-to-end chain beats a deep but partial one.
  \item \textbf{Every stage must be demonstrable.} If it cannot be shown, it is not done.
  \item \textbf{Honesty is a feature.} No claim of quantum advantage without a confidence interval and, eventually, a significance test. A null or negative result, measured rigorously, is a valid and reportable outcome.
\end{enumerate}

\subsection{The IBM Quantum hardware plan (already decided, not yet executed)}
\label{sec:ibm-plan}

The team has IBM Quantum cloud access (\code{qiskit-ibm-runtime==0.44.0} is already installed). The plan, fixed in \code{MVP\_SPEC.md} \S8.5 and unchanged since:

\begin{itemize}[leftmargin=1.4em]
  \item The \emph{full} classical-vs-quantum benchmark stays on the simulator, permanently. The full training kernel is $\approx$43{,}900 circuit evaluations $\approx$3--4 hours of QPU time, against a free-tier allowance on the order of 10 minutes/month --- roughly 20--40$\times$ over budget for a single run, before even considering repeated CV or multiple seeds.
  \item A \textbf{separate, small, scoped hardware validation experiment} is planned for Phase 6: a fixed 30-train/10-test subsample ($\approx$735 circuits, $\approx$3--4 minutes of QPU time), with the identical small kernel computed three ways --- exact statevector simulator, a noisy simulator, and the real IBM device --- then compared via element-wise kernel divergence, diagonal-fidelity deviation (a ground-truth-free noise indicator, since $K(x,x)$ should be exactly 1), and kernel-concentration statistics.
  \item This hardware run is executed \emph{once}, cached, and never queued live during a demo (queue times are unbounded).
  \item Hardware results are auditable (device name, calibration timestamp, job IDs, QPU seconds) but explicitly \emph{not} claimed to be bit-reproducible, because device calibration drifts daily.
\end{itemize}

% =========================================================================
\section{Engineering Principles Enforced Throughout}
\label{sec:principles}
% =========================================================================

These principles were established early and have been mechanically enforced (via code structure and tests, not just stated) in every phase since. A reader picking up Phase 6 should preserve all of them.

\subsection{Leakage safety}
\begin{itemize}[leftmargin=1.4em]
  \item The train/test split happens \emph{before} any statistic is ever computed (Phase 2), and is never re-created by a later phase.
  \item Every fitted transformer (imputer, encoder, scaler, feature selector, PCA) is fit on the training partition only, and only ever \code{.transform()}-ed on test/validation data.
  \item Cross-validation refits the \emph{entire} preprocessing chain inside every fold (Phase 3), using sklearn's own \code{GridSearchCV} machinery rather than a hand-rolled loop.
  \item \textbf{Structural leakage guards}: tuning functions (\code{run\_grid\_search}, \code{run\_qsvm\_grid\_search}) have \emph{no test-set parameter in their signature at all} --- verified directly via \code{inspect.signature()} in tests, not just by convention.
  \item \textbf{Behavioral leakage guards}: dedicated tests perturb an extreme, synthetic value into the test set and confirm training-derived statistics (imputer medians, scaler means, PCA axes, CV fold scores) do not change.
  \item The locked test set is touched by exactly one function, exactly once, after model selection is complete.
\end{itemize}

\subsection{Reproducibility}
\begin{itemize}[leftmargin=1.4em]
  \item A fixed random seed (42) propagates through every stochastic component: the split, cross-validation folds, \code{mutual\_info\_classif}'s internal estimator, every classical model's \code{random\_state}, bootstrap resampling.
  \item A deterministic \code{split\_id} (a short SHA-256-derived fingerprint of the split parameters) is computed once in Phase 2 and reproduced --- never re-derived --- by every subsequent phase. Its value for the default configuration is \code{e471025b07519a64}, and it appears verbatim in every phase's results and documentation.
  \item Every experiment's run record captures library versions, configuration, seeds, and timings as JSON.
  \item Immutable (\code{frozen=True}) configuration dataclasses (\code{PreprocessingConfig}, \code{QuantumConfig}) prevent one experiment's config from silently mutating and contaminating another's.
\end{itemize}

\subsection{Honesty / anti-overclaiming}
\begin{itemize}[leftmargin=1.4em]
  \item No document claims quantum advantage, quantum superiority, or clinical validity anywhere.
  \item Every results document distinguishes \emph{CV variability} (mean/std/t-distribution CI across folds, computed on training data only) from \emph{final test uncertainty} (percentile bootstrap CI on the single locked-test evaluation) and never conflates the two.
  \item Model-selection criterion is fixed \emph{before} results are observed (mean CV ROC-AUC, not accuracy) and is never changed retroactively to produce a better-looking result, even when a selected hyperparameter produces an ugly, degenerate outcome (this happened in Phase 5 --- see \S\ref{sec:phase5}).
  \item Every phase's documentation has an explicit ``What this does NOT establish'' section.
  \item Raw data immutability is checksum-verified (SHA-256) after every phase's code runs, not just claimed.
\end{itemize}

\subsection{Reuse, not duplication}
\begin{itemize}[leftmargin=1.4em]
  \item Phase 3 (classical) reuses Phase 2's preprocessing classes directly via thin sklearn-compatible adapter classes (\code{SharedFeatureTransformer}, \code{QuantumReadyTransformer}) rather than reimplementing preprocessing logic.
  \item Phase 4 (quantum infrastructure) reuses Phase 2's \code{run\_preprocessing\_pipeline()} directly for its data contract --- no new split, no new PCA fit.
  \item Phase 5 (QSVM) reuses Phase 3's scoring dictionary (\code{CV\_SCORING}), model-selection constant (\code{MODEL\_SELECTION\_METRIC}), and --- critically --- Phase 3's \emph{exact} \code{evaluate\_on\_test()} function for the final locked-test evaluation, unmodified. The same \code{ModelResult} dataclass is used for both classical and quantum models.
  \item Only one package (\code{src/quantum/}) is allowed to import Qiskit, so an API-breaking Qiskit upgrade has a confined blast radius.
\end{itemize}

\subsection{Documentation-in-sync-with-code}
Every phase produced a Markdown document in \code{docs/} that is explicitly declared to be \emph{derived from} the code and results, with a stated rule: if a number in the document ever disagrees with a fresh run of the code, the code's output is authoritative and the document must be corrected. Every non-trivial claim is paired with the specific test that verifies it (a ``Traceability'' table appears at the end of every phase document).

% =========================================================================
\section{Phase 1 --- Dataset Understanding and Validation}
\label{sec:phase1}
% =========================================================================

\phasestatus{\done}{\code{src/data/inspect\_dataset.py}}{\code{docs/DATASET.md}}{\code{tests/test\_dataset.py} (27 tests)}{---}

\subsection{What was built}
A read-only inspection and structural-validation module for \code{data/raw/processed.cleveland.data}. It loads the raw, headerless, comma-separated UCI file, assigns the 14 documented column names, treats \code{"?"} as the missing-value marker, and produces a full data-quality report: shape, dtypes, missing-value counts, duplicate detection, numeric/categorical profiles, domain-violation checks against the documented UCI value ranges, class distribution, and zero-variance-column detection. \textbf{No imputation, encoding, scaling, or target binarization happens in Phase 1} --- it is inspection only.

\subsection{Key verified facts about the dataset}

\begin{tabularx}{\textwidth}{@{}p{4.4cm}Y@{}}
\toprule
\textbf{Property} & \textbf{Value} \\
\midrule
Source file & \code{data/raw/processed.cleveland.data} \\
Rows / Columns & 303 / 14 (13 predictors + target \code{num}) \\
SHA-256 checksum & \code{\footnotesize a74b7efa387bc9d108d7d0115d831fe9b414b29ae7124f331b622b4efa0427c8} \\
Missing values & \code{ca}: 4 cells; \code{thal}: 2 cells; all others: 0 \\
Duplicate rows & 0 \\
Raw target \code{num} distribution & 0$\to$164, 1$\to$55, 2$\to$36, 3$\to$35, 4$\to$13 \\
Binary target preview (\code{num}$>$0) & 0 (no disease): 164 (54.1\%); 1 (at risk): 139 (45.9\%) \\
Categorical domain violations & 0 (every observed value matches the documented UCI codes) \\
\bottomrule
\end{tabularx}

\subsection{The 13 predictors and their semantic types}
\begin{tabularx}{\textwidth}{@{}p{3.5cm}p{3.6cm}Y@{}}
\toprule
\textbf{Group} & \textbf{Features} & \textbf{Notes} \\
\midrule
Continuous numeric & age, trestbps, chol, thalach, oldpeak & Median-imputed later, then scaled \\
Discrete numeric (ordinal count) & ca (0--3 vessels) & Kept numeric, not one-hot, to preserve order \\
Ordinal categorical & slope (3 levels) & Mode-imputed, kept as an ordered code \\
Binary categorical & sex, fbs, exang & Already 0/1; imputed only, never encoded/scaled \\
Nominal categorical & cp, restecg, thal & One-hot encoded after imputation \\
\bottomrule
\end{tabularx}

\subsection{Important findings worth remembering}
\begin{itemize}[leftmargin=1.4em]
  \item \textbf{The UCI documentation itself is inconsistent}: \code{heart-disease.names} states missing values are marked \code{-9.0}, but the actual processed file uses \code{"?"}. This was verified empirically (\code{-9.0} never appears; \code{"?"} appears exactly 6 times) and the loader trusts the empirical evidence, not the stale documentation.
  \item \textbf{\code{chol}=0 is a known undeclared missing-value sentinel at other UCI sites}, not Cleveland: verified that \code{chol}=0 occurs in 100\% of the 123 Switzerland records and 24.5\% of the 200 VA records, but \emph{zero} times in Cleveland. This directly motivated Phase 2's decision to declare sentinel codes per-column, per-dataset, in configuration rather than hardcoding them.
  \item \textbf{The clinical-leakage concern about \code{ca}/\code{thal}}: both are results of specialist cardiac investigations (fluoroscopy, thallium stress test) typically ordered once disease is already suspected. \code{ca} in particular sits close to a direct read-out of the angiographic diagnosis. This was flagged in Phase 1 and became the basis for Phase 2/3's ``screening'' feature-set ablation (features minus \code{ca}/\code{thal}), which Phase 3 confirmed empirically matters (see \S\ref{sec:phase3}).
  \item The inspection script was verified to \emph{generalize} correctly: running it against the Switzerland file independently rediscovered the \code{chol}=0 pattern and a 6.5\% minority-class imbalance warning, proving the logic is genuinely dataset-driven, not overfit to Cleveland's specific cleanliness.
\end{itemize}

% =========================================================================
\section{Phase 2 --- Leakage-Proof Preprocessing \& Feature Pipeline}
\label{sec:phase2}
% =========================================================================

\phasestatus{\done}{\code{src/preprocessing/\{config,pipeline,validation\}.py}}{\code{docs/PREPROCESSING.md}}{\code{tests/test\_preprocessing.py} (41 tests)}{---}

\subsection{What was built}
A pipeline that turns the raw dataset into \textbf{two} feature representations sharing \textbf{one} upstream stage, so the two branches can never silently diverge in how the data was cleaned:

\begin{center}
\begin{lstlisting}
RAW DATA
  -> TARGET CREATION (num -> target; target = 1 if num > 0 else 0; num preserved)
  -> STRATIFIED TRAIN/TEST SPLIT (80/20, seed 42)   <- split BEFORE any fitting
  -> PREPROCESSING (impute -> encode -> scale)       [fit on TRAIN only]
  -> FEATURE SELECTION (SelectKBest, mutual_info)    [fit on TRAIN only]
       |--------------------------|
       v                          v
  CLASSICAL-READY           PCA (4 comps) + range-normalize to [0,pi]  [fit on TRAIN only]
  FEATURES (10 cols)                |
                                     v
                            QUANTUM-READY FEATURES (4 cols, in [0, pi])
\end{lstlisting}
\end{center}

\subsection{Key configuration (\code{PreprocessingConfig}, defaults)}
\begin{tabularx}{\textwidth}{@{}p{4.4cm}Y@{}}
\toprule
\textbf{Field} & \textbf{Default} \\
\midrule
\code{random\_seed} & 42 \\
\code{test\_size} / \code{stratify} & 0.20 / True \\
\code{numeric\_impute\_strategy} & median \\
\code{categorical\_impute\_strategy} & most\_frequent \\
\code{scaler} & standard (also supports minmax, robust) \\
\code{feature\_selection\_method} & mutual\_info (also supports f\_classif) \\
\code{feature\_selection\_k} & 10 \\
\code{pca\_n\_components} & \textbf{4} --- this fixes the quantum branch's qubit count for all later phases \\
\code{quantum\_range} & $[0, \pi]$ \\
\code{feature\_set} & full (also supports ``screening'': excludes ca, thal) \\
\bottomrule
\end{tabularx}

\subsection{The locked split (referenced by every later phase)}
\begin{tabularx}{\textwidth}{@{}p{4.4cm}Y@{}}
\toprule
\textbf{Property} & \textbf{Value} \\
\midrule
Train / Test rows & 242 / 61 \\
\textbf{split\_id} & \textbf{\code{e471025b07519a64}} \\
Train class balance & 0: 131 (54.1\%), 1: 111 (45.9\%) \\
Test class balance & 0: 33 (54.1\%), 1: 28 (45.9\%) \\
Features after impute+encode+scale (pre-selection) & 17 (full feature set) \\
Selected features (top-10 by mutual information) & chol, thalach, ca, slope, sex, exang, cp\_2.0, cp\_4.0, thal\_6.0, thal\_7.0 \\
PCA cumulative explained variance (4 components) & 82.0\% \\
\bottomrule
\end{tabularx}

\subsection{Design decisions worth remembering}
\begin{itemize}[leftmargin=1.4em]
  \item \textbf{\code{MinMaxScaler(clip=True)} for the quantum range.} Discovered during development: a test-set PCA value can legitimately fall slightly outside the training-derived $[0,\pi]$ range (expected, since the scaler is fit on train only). \code{clip=True} guarantees the quantum-ready output is always bounded, without ever using test statistics to derive the bound.
  \item \textbf{\code{SharedFeaturePipeline}} is a hand-written two-stage object (not a bare sklearn \code{Pipeline}), because \code{feature\_selection\_k} must be clipped against the \emph{actual} post-one-hot-encoding column count, which is only known after the encoder has been fit.
  \item \textbf{The screening ablation} (\code{feature\_set="screening"}) is fully wired and tested but was \emph{not} run as an experiment in Phase 2 --- only prepared, per the phase's explicit scope. Phase 3 later ran it.
  \item \textbf{The dedicated leakage test} constructs a synthetic extreme value (cholesterol = 999{,}999) in the test set and proves the training-fitted scaler mean, imputer median, and PCA axes are bit-identical whether or not that value exists, because the test set is never part of any \code{.fit()} call.
\end{itemize}

% =========================================================================
\section{Phase 3 --- Classical ML Baseline \& Leakage-Safe Evaluation}
\label{sec:phase3}
% =========================================================================

\phasestatus{\done}{\code{src/classical/\{models,tuning,evaluation,train\_baselines\}.py}}{\code{docs/CLASSICAL\_BASELINE.md}}{\code{tests/test\_classical.py} (20 tests)}{\code{results/classical/**}}

\subsection{Purpose}
Establish a rigorous, fairly-tuned classical benchmark that later quantum models must be compared against. \textbf{Explicitly not} an attempt to show classical ML is weak.

\subsection{The four required models}
\begin{tabularx}{\textwidth}{@{}p{2.3cm}p{4.6cm}Y@{}}
\toprule
\textbf{Model} & \textbf{Library} & \textbf{Tuned parameter(s)} \\
\midrule
Logistic Regression & {\footnotesize\texttt{sklearn.linear\_\allowbreak model.\allowbreak LogisticRegression}} & C $\in \{0.001,\ldots,100\}$ (6 values) \\
RBF-Kernel SVM & {\footnotesize\texttt{sklearn.svm.\allowbreak SVC(kernel='rbf')}} & C (4), gamma (4) --- 16 combos \\
Random Forest & {\footnotesize\texttt{sklearn.\allowbreak ensemble.\allowbreak RandomForestClassifier}} & n\_estimators, max\_depth, min\_samples\_split, max\_features (24 combos) \\
XGBoost & {\footnotesize\texttt{xgboost.\allowbreak XGBClassifier}} (v3.2.0) & n\_estimators, max\_depth, learning\_rate, subsample, colsample (48 combos) \\
\bottomrule
\end{tabularx}

\textbf{A real sklearn-version bug was found and fixed, not hidden.} Installed \code{scikit-learn==1.8.0} deprecates the \code{penalty} keyword on \code{LogisticRegression} in favor of a unified \code{l1\_ratio} API; combining \code{penalty='l1'} with \code{solver='liblinear'} emitted an ``inconsistent values'' warning suggesting the l1 request might not be honored correctly. The fix: tune \code{C} only, with the standard L2 penalty via the \code{lbfgs} solver --- documented in code and in \code{CLASSICAL\_BASELINE.md} \S7.2, not silently patched over.

\subsection{Leakage-safe cross-validation architecture}
Two thin sklearn-compatible adapter classes wrap Phase 2's preprocessing so it can be dropped into \code{GridSearchCV}:
\begin{itemize}[leftmargin=1.4em]
  \item \code{SharedFeatureTransformer} --- delegates to Phase 2's \code{SharedFeaturePipeline}.
  \item \code{QuantumReadyTransformer} --- delegates to Phase 2's \code{build\_quantum\_pipeline} (PCA + range-normalize).
\end{itemize}
Both are placed as the first step(s) of an sklearn \code{Pipeline} together with the classifier, and the \emph{whole pipeline} (preprocessing included) is handed to \code{GridSearchCV(cv=StratifiedKFold(5, shuffle=True, seed=42))}. sklearn's own cross-validation machinery refits preprocessing fresh inside every fold. Model selection uses \textbf{mean CV ROC-AUC} (\code{MODEL\_SELECTION\_METRIC = "roc\_auc"}), never accuracy.

\subsection{The three experiments and their real results}

\textbf{Experiment 1 --- full feature set, classical-ready (10 selected features):}
\begin{center}
\small
\begin{tabular}{@{}lcccc@{}}
\toprule
\textbf{Model} & \textbf{CV ROC-AUC} & \textbf{Test ROC-AUC} & \textbf{Test Sens.} & \textbf{Test Spec.} \\
\midrule
Logistic Regression & 0.9019 $\pm$ 0.0130 & 0.9481 & 0.9286 & 0.8485 \\
RBF-Kernel SVM & 0.9055 $\pm$ 0.0189 & 0.9405 & 0.9286 & 0.8182 \\
Random Forest & 0.8894 $\pm$ 0.0193 & 0.9470 & 0.8214 & 0.9394 \\
XGBoost & 0.8871 $\pm$ 0.0216 & 0.9405 & 0.8571 & 0.8485 \\
\bottomrule
\end{tabular}
\end{center}
By the CV criterion, RBF-SVM ranks first (0.9055), narrowly ahead of Logistic Regression (0.9019) --- \emph{their 95\% CIs overlap substantially, this is not a decisive win}. On the single locked test set, Logistic Regression scores highest (0.9481). This CV-vs-test disagreement is reported honestly as the reason CV, not the single 61-row test score, is trusted as the primary model-selection evidence.

\textbf{Experiment 2 --- screening (11 features, \code{ca}/\code{thal} excluded), classical-ready:} Every single one of the four models' CV ROC-AUC \emph{dropped} when \code{ca}/\code{thal} were removed, consistently, by roughly 5 percentage points (mean across models: 0.896 $\to$ 0.845). This uniform effect across four architecturally different models is strong evidence that \code{ca}/\code{thal} genuinely carry substantial predictive signal beyond the other 11 features --- confirming the Phase 1 clinical-leakage concern was well-founded. \code{ca}/\code{thal} were never removed from the primary dataset; this is a config-only ablation.

\textbf{Experiment 3 --- full feature set, quantum-ready (4 PCA components):} This is the benchmark a future quantum model must be compared against fairly, because it uses the identical 4-dimensional representation a quantum model will see.
\begin{center}
\begin{tabular}{@{}lrr@{}}
\toprule
\textbf{Model} & \textbf{CV ROC-AUC (4-dim)} & \textbf{Test ROC-AUC} \\
\midrule
Logistic Regression & 0.8610 $\pm$ 0.0336 & 0.9210 \\
RBF-Kernel SVM & 0.8604 $\pm$ 0.0377 & 0.9242 \\
Random Forest & 0.8486 $\pm$ 0.0482 & 0.9167 \\
XGBoost & 0.8616 $\pm$ 0.0367 & 0.9275 \\
\midrule
\textbf{Mean} & \textbf{0.8579} & --- \\
\bottomrule
\end{tabular}
\end{center}
Compressing to 4 PCA components (82.0\% retained variance) costs $\approx$3.8 ROC-AUC points on average versus the full 10-feature representation --- a real, expected, disclosed cost of dimensionality reduction.

\textbf{Class-weighting side-check (Experiment 1's tuned hyperparameters, re-evaluated with \code{class\_weight='balanced'} / \code{scale\_pos\_weight}):} negligible effect on ROC-AUC (all four models changed by $<$0.002, well inside CV noise), with the expected small sensitivity-up/specificity-down trade-off. \textbf{Conclusion: class weighting is not used}, since the dataset is near-balanced (54.1\%/45.9\%) and weighting demonstrably does not help --- checked empirically, not assumed.

\subsection{Sanity check}
Every classical ROC-AUC lands well below 1.0 (0.887--0.906 CV, 0.94--0.95 test) --- plausible for this literature-standard dataset and \emph{not} suspiciously perfect, which is itself a basic sanity check against leakage.

% =========================================================================
\section{Phase 4 --- Quantum Data Pipeline \& Fidelity Kernel (Infrastructure)}
\label{sec:phase4}
% =========================================================================

\phasestatus{\done}{\code{src/quantum/\{config,data\_contract,feature\_maps,backends,kernel,sanity\_experiment\}.py}}{\code{docs/QUANTUM\_PIPELINE.md}}{\code{tests/test\_quantum.py} (53 Phase-4 tests)}{\code{results/quantum/sanity/}}

\subsection{Purpose and scope}
Build the infrastructure that turns an already-preprocessed patient record into a quantum kernel value --- deterministically, reproducibly, without touching the test set in any way that matters. \textbf{Explicitly not} a benchmark: no QSVM was trained, no accuracy/sensitivity/specificity was reported, nothing was executed on real IBM hardware.

\subsection{Critical environment decision: no \code{qiskit-aer}, no \code{qiskit-machine-learning}}
The installed \code{qiskit} version is \textbf{2.2.3} (a major, recent version). Installing \code{qiskit-aer} or \code{qiskit-machine-learning} would have pulled in \code{numpy>=2.0} --- a major version jump from the \code{numpy==1.26.4} that Phases 1--3's tested pipeline (pandas, scikit-learn, xgboost) runs on, verified via \code{pip install --dry-run}. \textbf{Decision: use only core Qiskit.} The exact statevector simulator, \code{qiskit.quantum\_info.Statevector.from\_instruction()}, needs no Aer at all --- at 4--6 qubits a statevector has 16--64 complex amplitudes, trivially cheap regardless of backend. This was verified by re-running the full Phase 1--3 test suite (68 tests) after adding Phase 4 code and confirming zero regressions and an unchanged \code{numpy} version.

\textbf{A real API-deprecation issue was found and worked around, not hidden.} \code{qiskit==2.2.3} deprecates the class-based \code{ZZFeatureMap}/\code{ZFeatureMap}/\code{PauliFeatureMap} (removal planned for Qiskit 3.0) in favor of function-based replacements: \code{zz\_feature\_map}, \code{z\_feature\_map}, \code{pauli\_feature\_map}. This implementation uses the function-based API throughout.

\subsection{The quantum data contract}
\code{load\_quantum\_dataset()} calls Phase 2's \code{run\_preprocessing\_pipeline()} directly --- the same function Phase 3's Experiment 3 used. \textbf{No new split is ever created.} Returns \code{X\_train\_quantum} (242$\times$4), \code{X\_test\_quantum} (61$\times$4), \code{y\_train}, \code{y\_test}, with \code{split\_id} reproduced (not re-derived) as \code{e471025b07519a64}. \code{n\_qubits} is derived from \code{PreprocessingConfig.pca\_n\_components} (never hardcoded), so the quantum pipeline can never silently disagree with Phase 2.

\subsection{The feature map}
\begin{tabularx}{\textwidth}{@{}p{4.4cm}Y@{}}
\toprule
\textbf{Property} & \textbf{Value} \\
\midrule
Feature map & \code{zz\_feature\_map} (default; \code{z\_feature\_map}, \code{pauli\_feature\_map} also supported) \\
Qubits & 4 (= \code{pca\_n\_components}) \\
Repetitions & 2 \\
Entanglement & linear (circular, full also supported) \\
Encoding & Angle encoding, one classical feature per qubit \\
Trainable quantum parameters & \textbf{Zero} --- the map is a fixed template; only the 4 data values are bound in \\
Free parameters (data) & 4 \\
Logical circuit depth (decomposed) & 19 \\
Gate counts (decomposed) & \{u: 22, cx: 12\} \\
\bottomrule
\end{tabularx}
\textbf{Why \code{zz\_feature\_map}}: entangling (\code{RZZ}-family) terms give the kernel access to pairwise feature correlations, unlike \code{z\_feature\_map} (no entanglement at all, essentially a product kernel). \textbf{Why linear entanglement}: fewest two-qubit gates, cheapest to simulate, and most compatible with a real device's limited connectivity (relevant for Phase 6, even though Phase 4 does not execute on hardware). \textbf{This configuration was deliberately not optimized or swept} --- a defensible standard starting point, not a validated best choice.

\subsection{Backend abstraction}
\begin{center}
\begin{lstlisting}
QuantumBackend (Protocol: compute_statevector(), capabilities())
  +-- StatevectorBackend      IMPLEMENTED  -- core Qiskit only, exact, deterministic
  +-- NoisySimulatorBackend   interface only -- raises NotImplementedError; a Phase 5/6 decision
  +-- IBMHardwareBackend      interface only -- raises NotImplementedError; scoped to Phase 6
\end{lstlisting}
\end{center}
Backend selection is purely config-driven (\code{QuantumConfig.backend\_name}); \code{kernel.py} never branches on the backend name itself.

\subsection{The fidelity kernel}
\begin{gather*}
K(x_i, x_j) = |\langle\phi(x_i)|\phi(x_j)\rangle|^2 \\
|\phi(x)\rangle = U_\phi(x)|0000\rangle
\end{gather*}
\begin{sloppypar}
Implemented efficiently: \code{compute\_statevectors(X, feature\_map, backend)} embeds every row \emph{once} (a statevector depends only on its own row), giving $O(n)$ circuit evaluations; \code{kernel\_matrix\_from\_statevectors(...)} assembles the matrix from cheap inner products, exploiting symmetry (upper triangle only) when applicable, but \textbf{the diagonal is always computed explicitly, never hardcoded to 1.0} --- a measured deviation from exactly 1.0 is itself a useful floating-point correctness signal.
\end{sloppypar}

\subsection{Sanity experiment (20 train / 10 test subsample, seed 42) --- real results}
\begin{tabularx}{\textwidth}{@{}p{4.4cm}Y@{}}
\toprule
\textbf{Item} & \textbf{Value} \\
\midrule
\code{K\_train\_train} / \code{K\_test\_train} / \code{K\_test\_test} shapes & (20,20) / (10,20) / (10,10) \\
Min / Max kernel value & 3.6e-05 / 1.0 \\
Diagonal deviation from 1.0 & 3.55e-15 (machine precision) \\
Reproducibility & Bit-identical recompute, max diff = 0.0 \\
Split ID & \code{e471025b07519a64} (matches Phase 2/3) \\
\bottomrule
\end{tabularx}
Explicitly labelled \code{"status": "sanity\_check\_only\_not\_a\_benchmark"} in its own JSON output.

% =========================================================================
\section{Phase 5 --- Full-Scale QSVM Experiment}
\label{sec:phase5}
% =========================================================================

\phasestatus{\done}{\code{src/quantum/\{qsvm,qsvm\_experiment\}.py}}{\code{docs/QUANTUM\_QSVM.md}}{21 new tests in \code{tests/test\_quantum.py} (74 total quantum tests)}{\code{results/quantum/phase5/}}

\subsection{Purpose}
Train the first real QSVM: \code{sklearn.svm.SVC(kernel='precomputed')} fed the Phase 4 fidelity kernel, computed at full scale on all 242 training and 61 test records, tuned leakage-safely, evaluated once on the locked test set, and compared descriptively (no significance test yet) against Phase 3's Experiment 3.

\subsection{Computational cost --- measured before committing to the design}
\begin{tabularx}{\textwidth}{Y r}
\toprule
\textbf{Step} & \textbf{Measured time} \\
\midrule
242 training-row statevectors & 0.449 s \\
61 test-row statevectors & 0.085 s \\
\code{K\_train\_train} assembly (242$\times$242, 29{,}161 pairs) & 0.091 s \\
\code{K\_test\_train} assembly (61$\times$242, 14{,}762 pairs) & 0.036 s \\
\textbf{Total kernel computation} & \textbf{0.662 s} \\
\code{GridSearchCV} (5 C values $\times$ 5 folds = 25 fits) & $\approx$0.4 s \\
\bottomrule
\end{tabularx}
This measurement is exactly why fine-grained mid-computation checkpointing was \emph{not} built: restartability is implemented at the level of caching the two final matrices to disk (keyed by feature-map config + \code{split\_id}), which is more than sufficient given sub-second real cost. Train and test statevectors are each computed \emph{exactly once} and reused for both matrices --- \code{compute\_kernel\_matrix\_cached()} from Phase 4 was deliberately \emph{not} used naively (it would have recomputed the 242 training statevectors a second time).

\subsection{A key methodological question, resolved and verified}
Is it leakage-safe to tune \code{C} by slicing the \emph{one} precomputed 242$\times$242 kernel per CV fold, rather than recomputing a fold-specific kernel? \textbf{Yes} --- because $K(x_i,x_j)$ is a fixed function of two feature vectors with zero fitted/learned state (the feature map is a template, not a model). This was \emph{empirically verified}, not just argued: a test compares \code{sklearn.model\_selection.cross\_val\_score} on the full precomputed kernel against a hand-written loop that manually slices \code{K[train\_idx][:,train\_idx]} / \code{K[val\_idx][:,train\_idx]} per fold --- the scores match exactly.

\textbf{One disclosed, bounded simplification}: the 4-dimensional PCA representation was fit \emph{once} on all 242 training rows (the single, locked Phase 2/4 contract), not refit per CV fold the way Phase 3's classical preprocessing was. This is a mild, inherited, explicitly-disclosed simplification upstream of Phase 5 itself, expected to have negligible practical impact at this sample size, and consistent with Phase 5's explicit instruction to reuse the locked representation unchanged.

\subsection{Hyperparameter tuning --- full CV grid table (real run)}
\begin{center}
\small
\begin{tabular}{@{}rrrrrr@{}}
\toprule
\textbf{C} & \textbf{CV ROC-AUC} & \textbf{CV Accuracy} & \textbf{CV Sensitivity} & \textbf{CV Specificity} & \textbf{Rank} \\
\midrule
\textbf{0.01} & \textbf{0.6945 $\pm$ 0.0281} & 0.5413 & \textbf{0.0000} & \textbf{1.0000} & \textbf{1 (selected)} \\
0.1 & 0.6945 $\pm$ 0.0285 & 0.5413 & 0.0000 & 1.0000 & 2 \\
1.0 & 0.6942 $\pm$ 0.0306 & 0.6403 & 0.5755 & 0.6943 & 3 \\
10.0 & 0.6629 $\pm$ 0.0737 & 0.6403 & 0.6032 & 0.6709 & 4 \\
100.0 & 0.6461 $\pm$ 0.0949 & 0.6406 & 0.5858 & 0.6869 & 5 \\
\bottomrule
\end{tabular}
\end{center}

\textbf{A finding reported precisely, not smoothed over}: the selected \code{C=0.01} has the (barely) highest CV ROC-AUC, tied within noise against C=0.1 and C=1.0, but \textbf{CV sensitivity = 0.0 and specificity = 1.0 at that C} --- the model predicts the majority class for every validation patient, in every fold. ROC-AUC (threshold-independent) can still register moderately above chance from ranking quality alone, even while the 0.5-threshold prediction has collapsed to one class. \textbf{This was not ``fixed'' by picking a different C} --- the selection criterion was fixed \emph{before} this result was observed (inherited from Phase 3 for consistency), and switching it post-hoc would itself be a form of outcome-chasing.

\subsection{Final locked-test evaluation (touched exactly once)}
\begin{center}
\begin{tabular}{lrl}
\toprule
\textbf{Metric} & \textbf{Value} & \textbf{95\% Bootstrap CI} \\
\midrule
Accuracy & 0.6721 & [0.557, 0.787] \\
\textbf{Sensitivity} & \textbf{0.7500} & [0.594, 0.909] \\
\textbf{Specificity} & \textbf{0.6061} & [0.437, 0.774] \\
Precision / F1 & 0.6176 / 0.6774 & --- \\
\textbf{ROC-AUC} & \textbf{0.7798} & [0.661, 0.888] \\
PR-AUC & 0.6839 & [0.527, 0.875] \\
\bottomrule
\end{tabular}
\end{center}
Confusion matrix: TN=20, FP=13, FN=7, TP=21 (n=61).

\textbf{A second finding reported plainly}: the final model, refit on all 242 rows, does \emph{not} show the CV folds' degenerate all-negative behavior --- sensitivity is a substantive 0.75. This CV-vs-final-refit discrepancy is documented as an open, unexplained observation (a plausible but \emph{unconfirmed} read: more training data shifts which points become support vectors), flagged explicitly for Phase 6 investigation rather than explained away.

\subsection{Descriptive comparison with Phase 3 (no significance test performed, by design)}
\begin{center}
\begin{tabular}{lrrrr}
\toprule
\textbf{Model} & \textbf{Accuracy} & \textbf{Sensitivity} & \textbf{Specificity} & \textbf{ROC-AUC} \\
\midrule
Logistic Regression & 0.869 & 0.857 & 0.879 & 0.921 \\
RBF-SVM & 0.869 & 0.893 & 0.848 & 0.924 \\
Random Forest & 0.869 & 0.857 & 0.879 & 0.917 \\
XGBoost & 0.836 & 0.857 & 0.818 & 0.928 \\
\textbf{QSVM (fidelity kernel)} & \textbf{0.672} & \textbf{0.750} & \textbf{0.606} & \textbf{0.780} \\
\bottomrule
\end{tabular}
\end{center}
All five rows use the identical 4-dimensional representation and identical locked split. \textbf{On every metric measured, the classical models outperform the QSVM on this dataset, at this qubit count, with this feature map.} Stated as a direct numeric fact. \textbf{What this does and does not establish} (verbatim from \code{QUANTUM\_QSVM.md} \S10):

\emph{Establishes:} the full pipeline works correctly end to end at scale; leakage-safe C-tuning on a precomputed quantum kernel is sound and verified; on this specific configuration, classical wins numerically.

\emph{Does NOT establish:} that quantum kernels are inferior in general, or under a different feature map/qubit count; that the CV-tie-but-different-sensitivity pattern is understood; that the CV-vs-final-test discrepancy has a confirmed cause; anything about statistical significance (not tested); anything about real hardware (pure simulation).

% =========================================================================
\section{Current Repository Structure}
% =========================================================================

\begin{lstlisting}
QuantumDx/
|-- PRD.md                          # Product Requirements Document (43 sections, never modified)
|-- MVP_SPEC.md                     # MVP spec + 13-phase plan (never modified)
|-- conftest.py                     # puts repo root on sys.path for pytest
|-- configs/
|   |-- preprocessing.yaml          # Phase 2 config (mirrors PreprocessingConfig)
|   `-- quantum.yaml                # Phase 4 config (mirrors QuantumConfig)
|-- data/
|   `-- raw/                        # UCI Heart Disease files; NEVER modified; checksum-verified
|-- src/
|   |-- data/
|   |   `-- inspect_dataset.py      # Phase 1
|   |-- preprocessing/
|   |   |-- config.py               # Phase 2
|   |   |-- pipeline.py             # Phase 2 (SharedFeaturePipeline, run_preprocessing_pipeline)
|   |   `-- validation.py           # Phase 2
|   |-- classical/
|   |   |-- models.py               # Phase 3 (4-model registry)
|   |   |-- tuning.py               # Phase 3 (sklearn adapters, CV_SCORING, MODEL_SELECTION_METRIC)
|   |   |-- evaluation.py           # Phase 3 (metrics, CIs, ModelResult, plots)
|   |   `-- train_baselines.py      # Phase 3 orchestrator
|   `-- quantum/                    # ONLY package that imports Qiskit
|       |-- config.py               # Phase 4 (QuantumConfig)
|       |-- data_contract.py        # Phase 4 (load_quantum_dataset)
|       |-- feature_maps.py         # Phase 4 (zz/z/pauli_feature_map)
|       |-- backends.py             # Phase 4 (StatevectorBackend + stubs)
|       |-- kernel.py               # Phase 4 (fidelity kernel + diagnostics + cache)
|       |-- sanity_experiment.py    # Phase 4 orchestrator (small subsample)
|       |-- qsvm.py                 # Phase 5 (SVC precomputed wrapper + CV)
|       `-- qsvm_experiment.py      # Phase 5 orchestrator (full 242/61 scale)
|-- tests/
|   |-- test_dataset.py             # 27 tests (Phase 1)
|   |-- test_preprocessing.py       # 41 tests (Phase 2)
|   |-- test_classical.py           # 20 tests (Phase 3)
|   `-- test_quantum.py             # 74 tests (Phase 4: 53, Phase 5: 21)
|-- results/
|   |-- classical/                  # Phase 3: 3 experiments, CSVs, JSON, 15 figures
|   `-- quantum/
|       |-- sanity/                 # Phase 4 sanity report + kernel matrices
|       `-- phase5/                 # Phase 5 kernels, CV grid, final metrics, comparison, figures
`-- docs/
    |-- DATASET.md                  # Phase 1
    |-- PREPROCESSING.md            # Phase 2
    |-- CLASSICAL_BASELINE.md       # Phase 3
    |-- QUANTUM_PIPELINE.md         # Phase 4
    |-- QUANTUM_QSVM.md             # Phase 5
    `-- latex/PROJECT_CONTEXT.tex   # THIS DOCUMENT
\end{lstlisting}

\textbf{Total test count as of Phase 5 completion: 162 tests, all passing, zero regressions across every phase.}

\subsection{Environment snapshot}
\begin{tabularx}{\textwidth}{@{}p{4.4cm}Y@{}}
\toprule
\textbf{Component} & \textbf{Version} \\
\midrule
Python & 3.11.9 (Windows-10-10.0.26200) \\
numpy & 1.26.4 (deliberately kept below 2.0 --- see Phase 4) \\
pandas & 2.3.3 \\
scikit-learn & 1.8.0 \\
scipy & 1.16.3 \\
xgboost & 3.2.0 \\
matplotlib & 3.10.1 \\
qiskit & 2.2.3 \\
qiskit-ibm-runtime & 0.44.0 (pre-installed, not yet wired up) \\
qiskit-aer / qiskit-machine-learning & \textbf{not installed, deliberately} \\
\bottomrule
\end{tabularx}

% =========================================================================
\section{How to Run Everything (command reference)}
% =========================================================================

\begin{lstlisting}
# Full test suite (all phases) -- ~35-60s, run this first to confirm the
# environment is healthy before doing anything else
python -m pytest -q

# Phase-by-phase test suites
python -m pytest tests/test_dataset.py -v        # Phase 1, 27 tests
python -m pytest tests/test_preprocessing.py -v  # Phase 2, 41 tests
python -m pytest tests/test_classical.py -v      # Phase 3, 20 tests
python -m pytest tests/test_quantum.py -v        # Phase 4+5, 74 tests

# Re-run the inspection / demonstration scripts (idempotent, safe)
python -m src.data.inspect_dataset
python -m src.preprocessing.pipeline
python -m src.classical.train_baselines      # ~9 minutes (3 experiments x 4 models)
python -m src.quantum.sanity_experiment       # <5 seconds
python -m src.quantum.qsvm_experiment         # <2 seconds (cache hit), ~1s cold
\end{lstlisting}

Before doing anything destructive, verify raw-data integrity:
\begin{lstlisting}
python -c "import hashlib; print(hashlib.sha256(open('data/raw/processed.cleveland.data','rb').read()).hexdigest())"
# must print: a74b7efa387bc9d108d7d0115d831fe9b414b29ae7124f331b622b4efa0427c8
\end{lstlisting}

% =========================================================================
\section{Future Phases}
% =========================================================================

\subsection{Phase 6 --- Scoped IBM Quantum Hardware Validation (next, not started)}
\pending

Per the plan fixed in \code{MVP\_SPEC.md} \S8.5 (\S\ref{sec:ibm-plan} above) and restated at the start of Phase 5:

\begin{enumerate}[leftmargin=1.4em]
  \item \textbf{Verify the IBM Quantum account/plan details first} --- do not assume: exact free-tier quota, session-vs-batch mode availability, which real device(s) are accessible, current queue times, and the pinned \code{qiskit-ibm-runtime} primitive API (\code{SamplerV2} interface) for the installed version (0.44.0).
  \item \textbf{Build \code{IBMHardwareBackend}} (currently a \code{NotImplementedError} stub in \code{src/quantum/backends.py}) implementing the same \code{QuantumBackend} Protocol used by \code{StatevectorBackend} --- so \code{kernel.py} needs zero changes.
  \item \textbf{Run the scoped validation experiment}: a fixed, seeded 30-train/10-test subsample of the same locked split ($\approx$735 circuits, $\approx$3--4 minutes of QPU time). Compute the identical kernel submatrix three ways: exact \code{StatevectorBackend}, a noisy simulator (may require finally installing \code{qiskit-aer} --- re-check the numpy 2.0 compatibility question first, since Phases 1--3 may tolerate it now, or may not; verify, do not assume), and the real IBM device.
  \item \textbf{Report}: element-wise kernel divergence (mean/max $|\Delta K|$) between simulator and hardware; diagonal-fidelity deviation on hardware (should be exactly 1, will not be --- a ground-truth-free noise indicator); kernel-concentration statistics per backend; job provenance (device name, calibration timestamp, job IDs, QPU seconds).
  \item \textbf{Do not} draw an accuracy/performance conclusion from the 10-point test subsample --- explicitly a divergence measurement, not a benchmark.
  \item \textbf{Never queue live during a demo.} Run once, cache to disk, serve from cache thereafter (mirrors the restartable-cache pattern already established in Phase 5).
  \item Credentials via a git-ignored \code{.env} (with a committed \code{.env.example}); never commit secrets.
\end{enumerate}

\textbf{Two concrete, open questions from Phase 5 worth investigating alongside or before Phase 6:}
\begin{itemize}[leftmargin=1.4em]
  \item Why did \code{C=0.01}/\code{0.1} produce a degenerate (sensitivity=0) classifier in every CV fold, yet the full-242-row refit at the same C did not? (Phase 5, \S\ref{sec:phase5})
  \item Is the QSVM's low off-diagonal kernel mean ($\approx$0.09 in \code{K\_train\_train}) an early sign of kernel concentration, and would a feature-map sweep (different \code{reps}, entanglement, or \code{pauli\_feature\_map}) change the classical-vs-quantum gap observed in Phase 5?
\end{itemize}

\subsection{Later phases (per \code{MVP\_SPEC.md}'s 13-phase plan, not yet started)}
\pending

\begin{itemize}[leftmargin=1.4em]
  \item \textbf{VQC/QNN} --- the second quantum model family named in SIH26139. Explicitly the \emph{one pre-authorized droppable item} in the MVP plan if time runs short; QSVM alone already satisfies the problem statement's ``such as QSVM, QNN, or VQC'' language.
  \item \textbf{Statistical significance testing} (McNemar for accuracy, DeLong or paired bootstrap for AUC) comparing QSVM vs. the Phase 3 classical baselines properly --- explicitly deferred out of Phase 5 by instruction, and the natural next step once Phase 6's hardware context exists.
  \item \textbf{Explainability} --- SHAP for classical models; permutation importance + KernelSHAP (model-agnostic) for the QSVM, since quantum models expose no native feature attributions; a PCA loading map translating the 4 quantum components back to clinical features (already computed once in Phase 2, not yet surfaced in a dedicated explainability module).
  \item \textbf{Robustness / generalization suite} --- repeated CV variance (already partially done via bootstrap CIs), noise-injection degradation curves, subgroup fairness slices (by sex, age band), and a cross-site check against the Hungarian/Switzerland/VA UCI files (available in \code{data/raw/} but never used for training, only for the Phase 1 generalization check of the inspection script itself).
  \item \textbf{FastAPI backend} --- 9 endpoints as scoped in \code{MVP\_SPEC.md} \S14, deliberately excluding a \code{POST /train} endpoint (training stays CLI-driven; the API only serves stored results and predictions).
  \item \textbf{Dashboard} --- React (or a Streamlit fallback, pre-authorized if time is short) with the 7 pages specified in \code{PRD.md}: Home, Dataset Analysis, Patient Prediction, Model Comparison, Explainability, Quantum Circuit/Model Info, Experiment Results.
  \item \textbf{Final SIH demo preparation} --- code freeze, a rehearsed 10-minute run sheet (already drafted in \code{MVP\_SPEC.md} \S25.2), an offline fallback recording, and the acceptance checklist in \code{MVP\_SPEC.md} \S19.
\end{itemize}

% =========================================================================
\section{Key Numbers Quick-Reference}
% =========================================================================

\begin{tabularx}{\textwidth}{@{}p{5.6cm}Y@{}}
\toprule
\textbf{Quantity} & \textbf{Value} \\
\midrule
Dataset & UCI Heart Disease, Cleveland, 303 rows, 14 columns \\
Raw file SHA-256 & \code{\footnotesize a74b7efa387bc9d108d7d0115d831fe9b414b29ae7124f331b622b4efa0427c8} \\
Locked split\_id & \code{e471025b07519a64} \\
Train / Test & 242 / 61 (seed 42, stratified 80/20) \\
Selected classical features & 10 (of 17 encoded, from 13 raw) \\
Quantum features / qubits & 4 (= \code{PreprocessingConfig.pca\_n\_components}) \\
Quantum encoding range & $[0, \pi]$ \\
Feature map & \code{zz\_feature\_map}, reps=2, linear entanglement \\
Best classical model (Exp.\ 1, CV) & RBF-SVM, ROC-AUC 0.9055 $\pm$ 0.019 \\
Best classical model (Exp.\ 1, test) & Logistic Regression, ROC-AUC 0.9481 \\
Classical quantum-ready baseline (Exp.\ 3, mean CV) & ROC-AUC 0.858 \\
QSVM (Phase 5, test) & ROC-AUC 0.7798 [CI 0.661, 0.888] \\
Total tests passing & 162 / 162 \\
Phases complete & 1, 2, 3, 4, 5 \\
Next phase & 6 (scoped IBM hardware validation) \\
\bottomrule
\end{tabularx}

\vspace{1em}
\noindent\textit{This document was generated after Phase 5's completion and review. It should be regenerated or amended after each future phase to remain a reliable handoff artifact.}

\end{document}
